\begin{appendices}
\renewcommand{\thesection}{A\arabic{section}}
\renewcommand{\thesubsection}{A\arabic{section}.\arabic{subsection}}
\renewcommand{\thefigure}{A\arabic{figure}}
\renewcommand{\thetable}{A\arabic{table}}

% -------------------------------------------------------------------
\section{Demostración Extendida del Teorema de \texorpdfstring{$\phi$}{phi}-Mixing}
\label{apx:teorema_phi_mixing}

Sea \( \{S_t\} \) una cadena de Markov homogénea con espacio de estados finito \( S = \{s_1, \dots, s_N\} \), matriz de transición \( P \), e irreducible y aperiódica.

\begin{theorem}[\cite{Bradley2005}]
Una cadena de Markov \( \{S_t\} \) es \( \phi \)-mixing con tasa geométrica si existe \( \delta > 0 \) tal que para todo par de estados \( s_i, s_j \in S \), existe \( t \in \mathbb{N} \) con:
\[ P^t(s_j | s_i) > \delta. \]
\end{theorem}

\begin{proof}[Idea de la demostración]
El resultado se deduce a partir del hecho de que las cadenas de Markov irreducibles, aperiódicas y con espacio de estados finito son uniformemente ergódicas. Tal propiedad implica mezcla fuerte \( \phi \)-mixing con decaimiento exponencial.
\end{proof}

\textbf{Observación práctica:} En el modelo TCROC-Markov, la irreducibilidad puede garantizarse si el rango de \( \alpha_{t,\lambda,W} \) permite transiciones entre todos los estados definidos por \( \pi \).

% -------------------------------------------------------------------
\section{Estabilidad Discreta de la Función de Cuantización
\texorpdfstring{$\pi$}{pi}}
\label{apx:lipschitz_constante_pi}

La función de cuantización $\pi\colon \mathbb{R} \to \mathcal{S}$
definida en la Definición~\ref{def:pi_cuantizacion} es una función
escalonada y, por tanto, discontinua en los umbrales
$\{\theta_j\}$.  En sentido estricto, $\pi$ no es Lipschitz
continua.  Sin embargo, la estabilidad del modelo se garantiza
mediante la separación entre umbrales, como se demuestra en el
Lema~\ref{lem:estabilidad_discreta}.

Si se desea una versión suavizada de $\pi$ para análisis
asintótico, se puede considerar la aproximación sigmoide:
\[
\pi^*(x) = \sum_{j=1}^{K} s_j \cdot h_j(x), \qquad
h_j(x) = \frac{1}{1 + \exp(-a(x - \theta_j))},
\]
donde $a > 0$ controla la pendiente de la transición.
La constante de Lipschitz de $\pi^*$ satisface
$L_{\pi^*} \leq a/4$, alcanzada en los umbrales ($x = \theta_j$).
Cuando $a \to \infty$, $\pi^* \to \pi$ puntualmente, pero
$L_{\pi^*} \to \infty$, lo cual es consistente con la
discontinuidad de $\pi$.  Esta suavización es útil para
demostraciones técnicas, pero en la práctica del modelo
TCROC-Markov se trabaja directamente con la versión discreta.

% -------------------------------------------------------------------
\section{Ilustraciones del Proceso Estocástico \texorpdfstring{$\{S_t\}$}{St}}
\label{apx:figuras_proceso_st}

\begin{figure}[htbp]
\centering
\begin{tikzpicture}
  \begin{axis}[
    width=0.8\textwidth,
    xlabel={Tiempo \( t \)}, ylabel={Estado \( S_t \)},
    ymin=0.5, ymax=4.5,
    ytick={1,2,3,4}, yticklabels={\textnormal{s\_1},\textnormal{s\_2},\textnormal{s\_3},\textnormal{s\_4}},
    legend pos=north west,
    grid=major,
    title={Trayectorias simuladas del proceso \( \{S_t\} \)}
  ]
  \addplot+[mark=*, thick, nordFrost] coordinates {
    (1,2) (2,2) (3,1) (4,1) (5,2) (6,3) (7,4) (8,3) (9,2) (10,1)
  };
  \addplot+[mark=triangle*, thick, nordRed] coordinates {
    (1,3) (2,3) (3,2) (4,2) (5,2) (6,2) (7,1) (8,1) (9,2) (10,3)
  };
  \legend{\( \lambda = 0.95 \), \( \lambda = 0.75 \)}
  \end{axis}
\end{tikzpicture}
\caption{Trayectorias simuladas del proceso \( \{S_t\} \) para diferentes valores de \( \lambda \): menor \( \lambda \) implica mayor variabilidad y menor persistencia.}
\label{fig:apx_st_trayectorias}
\end{figure}

\begin{figure}[htbp]
\centering
\begin{tikzpicture}[node distance=2cm and 2.5cm, >=Stealth]
\node[state, fill=nordRed!15] (s1) {\( s_1 \)};
\node[state, fill=nordFrost!15] (s2) [right=of s1] {\( s_2 \)};
\node[state, fill=nordGreen!15] (s3) [below=of s2] {\( s_3 \)};
\node[state, fill=nordYellow!15] (s4) [left=of s3] {\( s_4 \)};

\draw[->] (s1) to[bend left] node[above] {0.3} (s2);
\draw[->] (s2) to[bend left] node[below] {0.2} (s1);
\draw[->] (s2) to[bend left] node[right] {0.5} (s3);
\draw[->] (s3) to[bend left] node[left] {0.4} (s2);
\draw[->] (s3) to[bend left] node[below] {0.3} (s4);
\draw[->] (s4) to[bend left] node[above] {0.1} (s3);
\draw[->] (s1) to[loop above] node {0.4} (s1);
\draw[->] (s2) to[loop above] node {0.3} (s2);
\draw[->] (s3) to[loop below] node {0.2} (s3);
\draw[->] (s4) to[loop below] node {0.6} (s4);
\end{tikzpicture}
\caption{Diagrama simplificado de transiciones entre estados \( s_j \). Las probabilidades se estiman a partir de la matriz \( \hat{P} \).}
\label{fig:apx_diagrama_transiciones}
\end{figure}

\vspace{1em}
\noindent\textbf{Nota:} Las figuras de esta sección ilustran el comportamiento cualitativo del proceso $\{S_t\}$ y sus transiciones. Para las simulaciones basadas en las matrices $\hat{P}$ estimadas con datos reales, véase el Capítulo~\ref{cap:regimenes_combustibles}.




% -------------------------------------------------------------------
% -------------------------------------------------------------------
% FINAL APÉNDICE A: Profundización Técnica del Modelo TCROC-Markov
% -------------------------------------------------------------------
% -------------------------------------------------------------------
% -------------------------------------------------------------------
% -------------------------------------------------------------------
% APÉNDICE B: Código para Análisis de Sensibilidad
% -------------------------------------------------------------------

% --- Configuracion para contadores del Apendice B ---
\renewcommand{\thesection}{B\arabic{section}}
\renewcommand{\thesubsection}{B\arabic{section}.\arabic{subsection}}
\renewcommand{\thefigure}{B\arabic{figure}}
\renewcommand{\thetable}{B\arabic{table}}
\setcounter{section}{0} % Reinicia el contador para que la siguiente seccion sea B1
\setcounter{figure}{0}
\setcounter{table}{0}

\section{Codigo Python para Analisis de Sensibilidad}
\label{apx:codigo_sensibilidad}

A continuacion, se presenta el codigo completo en Python utilizado para realizar el analisis de sensibilidad de los parametros de discretizacion $W$ y $\lambda$, como se describe en la subseccion de "Analisis de Sensibilidad de los Parametros de Discretizacion".

\begin{lstlisting}[language=Python, caption={Script de Python para el calculo de distribucion de estados en el analisis de sensibilidad.}, label={lst:codigo_sensibilidad}]
import numpy as np
import pandas as pd

# --- Funciones Auxiliares ---

def calculate_alpha(v, W, lam):
    """Calcula la serie de tasas relativas alpha_t usando la formula de TCROC."""
    alphas = []
    T = len(v)
    weights = np.array([lam**(t) for t in range(W)])[::-1] 
    
    for t in range(W, T):
        v_current_window = v[t-W+1 : t+1]
        v_previous_window = v[t-W : t]
        
        numerator = np.sum(weights * v_current_window * v_previous_window)
        denominator = np.sum(weights * v_previous_window**2)
        
        if denominator == 0:
            alpha = 0
        else:
            alpha = -1 + (numerator / denominator)
        
        alphas.append(alpha)
        
    return np.array(alphas)

def assign_states(alphas):
    """Asigna estados discretos segun los umbrales definidos."""
    conditions = [
        alphas < -0.01,
        (alphas >= -0.01) & (alphas <= 0.02),
        (alphas > 0.02) & (alphas <= 0.04),
        alphas > 0.04
    ]
    choices = [1, 2, 3, 4]
    
    return np.select(conditions, choices, default=0)

# --- Script Principal del Analisis de Sensibilidad ---

# 1. DATOS Y PARAMETROS
datos_reales_pib = np.array([
    16214.3, 16211.7, 16381.9, 16692.9, 17296.5, 17903.1, 18634.1, 
    19337.8, 19708.9, 18813.3, 19104.7, 19431.7, 19828.9, 19983.0, 
    20199.1, 20579.2, 20982.7, 21595.6, 22019.3, 22177.7, 19838.3, 
    21961.6, 22491.3, 22900.9
])

ws_to_test = [2, 3, 4, 5]
lambdas_to_test = np.arange(1.0, 0.79, -0.01)
params_to_test = [(w, lam) for w in ws_to_test for lam in lambdas_to_test]

# 2. REALIZAR EL ANALISIS
results = []
for W, lam in params_to_test:
    alphas = calculate_alpha(datos_reales_pib, W, lam)
    states = assign_states(alphas)
    
    unique_states, counts = np.unique(states, return_counts=True)
    state_counts = dict(zip(unique_states, counts))
    
    total_states = len(states)
    if total_states > 0:
        s1_perc = (state_counts.get(1, 0) / total_states) * 100
        s2_perc = (state_counts.get(2, 0) / total_states) * 100
        s3_perc = (state_counts.get(3, 0) / total_states) * 100
        s4_perc = (state_counts.get(4, 0) / total_states) * 100
    else:
        s1_perc, s2_perc, s3_perc, s4_perc = 0, 0, 0, 0

    results.append({
        'W': W,
        'lambda': lam,
        '% s1': s1_perc,
        '% s2': s2_perc,
        '% s3': s3_perc,
        '% s4': s4_perc
    })

# 3. MOSTRAR RESULTADOS
df_results = pd.DataFrame(results)
pd.set_option('display.max_rows', None)
print(df_results.round(2))
\end{lstlisting}

% -------------------------------------------------------------------
% APÉNDICE B: Código para Análisis de Robustez Predictiva
% -------------------------------------------------------------------
\section{Código Python para Validación con Múltiples Particiones}
\label{apx:codigo_validacion_completo}

A continuación, se presenta el código completo en Python utilizado para realizar el análisis de robustez de la validación predictiva, como se resume en la Tabla~\ref{tab:rendimiento_predictivo_resumen}. El script itera sobre múltiples particiones de entrenamiento/prueba para evaluar la estabilidad del rendimiento del modelo.

\begin{lstlisting}[language=Python, caption={Script de Python para validación con múltiples particiones.}, label={lst:codigo_validacion_multiples_particiones}]
import numpy as np
import pandas as pd
from scipy.optimize import nnls

# --- Funciones Auxiliares ---

def calculate_alpha(v, W, lam):
    """Calcula la serie de tasas relativas alpha_t usando la formula de TCROC."""
    alphas = []
    T = len(v)
    weights = np.array([lam**(t) for t in range(W)])[::-1]
    for t in range(W, T):
        v_current_window = v[t-W+1 : t+1]
        v_previous_window = v[t-W : t]
        numerator = np.sum(weights * v_current_window * v_previous_window)
        denominator = np.sum(weights * v_previous_window**2)
        if denominator == 0:
            alpha = 0
        else:
            alpha = -1 + (numerator / denominator)
        alphas.append(alpha)
    return np.array(alphas)

def assign_states(alphas):
    """Asigna estados discretos segun los umbrales definidos."""
    conditions = [
        alphas < -0.01,
        (alphas >= -0.01) & (alphas <= 0.02),
        (alphas > 0.02) & (alphas <= 0.04),
        alphas > 0.04
    ]
    choices = [1, 2, 3, 4]
    return np.select(conditions, choices, default=0)

def SRep(datos, ss):
    """Tu funcion SRep original (del Listado 1.2, adaptada)."""
    n0 = datos.shape[0]
    if ss >= datos.shape[1]:
        ss = datos.shape[1] - 1
        
    S0_data = datos[:, :ss]
    S0 = np.kron(S0_data, np.identity(n0)).T
    S1 = S0.T @ (datos[:, 1:(1 + ss)].T).reshape(ss * n0, 1)
    
    C = np.kron(np.identity(n0), np.ones((1, n0)))
    
    Mr = np.zeros((n0**2 + n0, n0**2))
    Mr[:n0**2, :] = S0.T @ S0
    Mr[n0**2:, :] = C
    
    rhs = np.zeros((n0**2 + n0))
    rhs[:n0**2] = S1.flatten()
    rhs[n0**2:] = 1
    
    c, rnorm = nnls(Mr, rhs)
    
    Pr = c.reshape(n0, n0).T
    
    col_sums = Pr.sum(axis=0)
    non_zero_cols = col_sums > 0
    if np.any(non_zero_cols):
        Pr[:, non_zero_cols] = np.divide(Pr[:, non_zero_cols], col_sums[non_zero_cols], where=col_sums[non_zero_cols]!=0)

    return Pr

def calculate_count_matrix(states_sequence, num_states=4):
    """Calcula la matriz de conteo de transiciones C."""
    count_matrix = np.zeros((num_states, num_states), dtype=int)
    for i in range(len(states_sequence) - 1):
        from_state = states_sequence[i] - 1
        to_state = states_sequence[i+1] - 1
        if from_state >= 0 and to_state >= 0:
            count_matrix[to_state, from_state] += 1
    return count_matrix

# --- Script Principal con Multiples Divisiones ---

# 1. DATOS Y PARAMETROS
datos_reales_pib = np.array([
    16214.3, 16211.7, 16381.9, 16692.9, 17296.5, 17903.1, 18634.1, 
    19337.8, 19708.9, 18813.3, 19104.7, 19431.7, 19828.9, 19983.0, 
    20199.1, 20579.2, 20982.7, 21595.6, 22019.3, 22177.7, 19838.3, 
    21961.6, 22491.3, 22900.9
])
W_model = 2
lambda_model = 1.0
NUM_STATES = 4

# 2. GENERACION DE DATOS BASE
alphas_full = calculate_alpha(datos_reales_pib, W_model, lambda_model)
states_full = assign_states(alphas_full)
num_samples = len(states_full)
X_full = np.zeros((NUM_STATES, num_samples))
for i, state in enumerate(states_full):
    if state > 0:
        X_full[state - 1, i] = 1

C_full = calculate_count_matrix(states_full, num_states=NUM_STATES)
P_hat_full = SRep(X_full, ss=X_full.shape[1] - 1)

# 3. BUCLE PARA PROBAR DIFERENTES DIVISIONES
split_ratios_to_test = np.arange(0.60, 1.0, 0.05)

for ratio in split_ratios_to_test:
    train_perc = int(ratio * 100)
    test_perc = 100 - train_perc
    
    print("\n" + "#"*70)
    print(f"###   RESULTADOS PARA DIVISION {train_perc}% / {test_perc}%   ###")
    print("#"*70 + "\n")

    # DIVISION EN ENTRENAMIENTO Y PRUEBA
    split_index = int(num_samples * ratio)
    states_train = states_full[:split_index]
    states_test = states_full[split_index:]
    X_train = X_full[:, :split_index]

    # CALCULO DE MATRICES
    C_train = calculate_count_matrix(states_train, num_states=NUM_STATES)
    C_test = calculate_count_matrix(states_test, num_states=NUM_STATES)
    P_hat_trained = SRep(X_train, ss=X_train.shape[1] - 1)

    # EVALUACION PREDICTIVA
    predictions = []
    actual_outcomes = []
    for i in range(len(states_test) - 1):
        current_state_idx = states_test[i] - 1
        prob_next_state = P_hat_trained[:, current_state_idx]
        predicted_state = np.argmax(prob_next_state) + 1
        predictions.append(predicted_state)
        actual_next_state = states_test[i+1]
        actual_outcomes.append(actual_next_state)

    # SALIDA PARA ESTA DIVISION
    print("--- Matriz de Conteo C (Entrenamiento) ---")
    print(pd.DataFrame(C_train, index=[f'a s{i+1}' for i in range(4)], columns=[f'de s{i+1}' for i in range(4)]))
    
    print("\n--- Matriz de Conteo C (Prueba) ---")
    print(pd.DataFrame(C_test, index=[f'a s{i+1}' for i in range(4)], columns=[f'de s{i+1}' for i in range(4)]))

    print("\n--- Matriz de Conteo C (Completa) ---")
    print(pd.DataFrame(C_full, index=[f'a s{i+1}' for i in range(4)], columns=[f'de s{i+1}' for i in range(4)]))

    print("\n--- Matriz P_hat (Entrenamiento) ---")
    print(pd.DataFrame(P_hat_trained, index=[f'a s{i+1}' for i in range(4)], columns=[f'de s{i+1}' for i in range(4)]).round(4))

    print("\n--- Matriz P_hat (Completa) ---")
    print(pd.DataFrame(P_hat_full, index=[f'a s{i+1}' for i in range(4)], columns=[f'de s{i+1}' for i in range(4)]).round(4))
    print("\n" + "="*50 + "\n")

    print("--- Desglose de Predicciones en el Conjunto de Prueba ---")
    prediction_df = pd.DataFrame({
        'Estado Actual (t)': states_test[:-1],
        'Estado Real Siguiente (t+1)': actual_outcomes,
        'Estado Predicho (t+1)': predictions
    })
    if not prediction_df.empty:
        prediction_df['Correcto?'] = (prediction_df['Estado Real Siguiente (t+1)'] == prediction_df['Estado Predicho (t+1)'])
    print(prediction_df)

    # METRICAS RESUMEN
    correct_predictions = prediction_df['Correcto?'].sum() if 'Correcto?' in prediction_df else 0
    total_predictions = len(prediction_df)
    accuracy = correct_predictions / total_predictions if total_predictions > 0 else 0

    print("\n--- Metricas Resumen ---")
    print(f"Predicciones correctas: {correct_predictions} de {total_predictions}")
    print(f"Precision (Accuracy): {accuracy:.3f}")
\end{lstlisting}

% -------------------------------------------------------------------
\renewcommand{\thesection}{B\arabic{section}}
\setcounter{section}{0}
\renewcommand{\thesubsection}{B\arabic{section}.\arabic{subsection}}

\section{Código Fuente: Estimación y Simulación}
\label{apx:codigo_srep}

El siguiente listado contiene la implementación del algoritmo
\texttt{SRep} para la estimación de matrices de transición
mediante Mínimos Cuadrados No Negativos (NNLS), así como
la simulación de la evolución de la cadena de Markov.

\begin{listing}[H]
\begin{python}
import numpy as np
from scipy.linalg import kron, identity
from scipy.optimize import nnls

def SRep(datos, ss):
    n0 = datos.shape[0]
    S0 = datos[:, :ss]
    S0 = kron(S0, identity(n0)).T
    S1 = S0.T @ (datos[:, 1:(1 + ss)].T).reshape(ss * n0)
    C = kron(identity(n0), np.ones((1, n0)))
    Mr = np.zeros((n0**2 + n0, n0**2))
    Mr[:n0**2, :] = S0.T @ S0
    Mr[n0**2:, :] = C
    rhs = np.zeros((n0**2 + n0))
    rhs[:n0**2] = S1
    rhs[n0**2:] = 1
    c = np.zeros((n0**2, 1))
    c[:, 0] = nnls(Mr, rhs)[0]
    Pr = c.reshape(n0, n0).T
    Pr = Pr @ np.diag(1 / Pr.sum(axis=0))
    return Pr
\end{python}
\caption{Implementación del algoritmo \texttt{SRep} para estimar
matrices de transición estocásticas mediante NNLS.}
\label{lst:srep_implementacion}
\end{listing}

\begin{listing}[H]
\begin{python}
# Simulacion del proceso de Markov
T = 20
p0 = np.zeros((4, T))
p0[0, 0] = 1  # Estado inicial: s1

for j in range(T - 1):
    p0[:, j+1] = P_hat @ p0[:, j]
\end{python}
\caption{Simulación de la evolución temporal de una cadena de
Markov a partir de la matriz de transición $\hat{P}$.}
\label{lst:simulacion_markov}
\end{listing}

\end{appendices}

